\documentclass[aps,prf,reprint,twocolumn]{revtex4-2}
\usepackage{graphicx}
\usepackage{amsmath,amssymb}
\usepackage{hyperref}
\usepackage{booktabs}

\newcommand{\mec}{\text{MEC}}

\begin{document}

\title{Spectral Sustainment and Stable Broad-Spectrum Organization in Neural Dynamics}

\author{}

\date{\today}

\begin{abstract}
Systems maintaining high-dimensional dynamics under stabilization occupy a restricted spectral regime characterized by slow spectral decay and broad eigenvalue support. We analyze 40 recordings from rodent medial entorhinal cortex (MEC) during spatial navigation and 22 synthetic dynamical conditions. Spectral decay rate $\alpha \approx 0.067$ for MEC separates broad-spectrum organization ($\alpha < 0.1$) from dimensional collapse ($\alpha > 0.5$). The quantity $\alpha\cdot\text{PR} \approx 2.3$ serves as a conserved invariant distinguishing MEC from both collapsed and critically organized systems. Finite-size scaling analysis reveals that MEC/sparse ensemble separation grows with system size $N$ as $d(N) \sim N^{2.2}$, confirming that the observed broad-spectrum correlation geometry reflects genuine structure rather than finite-size artifacts. The dominant organization resides in instantaneous precision geometry ($\alpha \approx 0.33$, PR $\approx 98$, $78\%$ sparsity) rather than temporal propagation. We propose a candidate stable broad-spectrum regime---provisionally termed Stably Sustained Critical Spectrum (SSCS)---characterized by $\alpha < 0.1$, PR $> 15$, and LC $< -5$.
\end{abstract}

\maketitle

\section{Introduction}

Dynamical systems that sustain high-dimensional activity while maintaining stability are rare. Most stabilized systems collapse to low-dimensional manifolds characterized by rapid spectral decay and few active modes \cite{ganguli2012,pennington2017}. A small class of systems---including neural populations during active behavior \cite{mazor2005,churchland2012,stringer2019} and certain turbulent flows \cite{holmes2012}---maintain broad spectral support with many active degrees of freedom.

Understanding the spectral conditions that separate broad-spectrum sustainment from collapse is a statistical physics problem with implications for neural computation, complex systems, and the theory of high-dimensional organization.

Here we characterize the spectral sustainment regime observed in recordings from rodent medial entorhinal cortex. We identify a spectral decay parameter $\alpha$ that serves as an organizing variable, establish scaling relationships, compare against random matrix ensembles, and discuss the evidence for a candidate stable broad-spectrum regime.

\section{Results}

\subsection{Spectral organization of MEC dynamics}

The eigenvalue spectrum of the correlation matrix of MEC activity follows a slow exponential decay with exponent $\alpha = 0.067 \pm 0.03$ (Fig.~\ref{fig:spectra}, mean across 40 sessions). The exponential form $\lambda_k \sim e^{-\alpha k}$ is dictated by the spectral bound $\text{PR} \ge (1+e^{-\alpha})/(1-e^{-\alpha})$ derived from the geometric series of the covariance spectrum \cite{trefethen1997}. This decay is substantially slower than that of random matrix ensembles \cite{wigner1955,mehta2004} (GOE: $\alpha \approx 0.27$ for matched dimensionality) and of synthetic stabilized systems ($\alpha > 0.5$).

The participation ratio PR $= 34 \pm 12$ confirms broad spectral support, with approximately 30\% of dimensions carrying significant variance. The closure metric LC $= -3.5 \pm 1.2$ indicates moderate trajectory collapse, consistent with stable but non-degenerate dynamics.

\begin{figure}[t]
\centering
\includegraphics[width=\columnwidth]{figures/fig1_mec_spectra.pdf}
\caption{\textbf{Eigenvalue spectra of MEC dynamics.} (a) Mean eigenvalue decay (black) with $\pm 1\sigma$ confidence interval (shaded) across 40 MEC recordings. GOE prediction for matched $N$ (red dashed) and synthetic stabilized system (orange dashed) shown for comparison. (b) Distribution of $\alpha$ across recordings.}
\label{fig:spectra}
\end{figure}

\subsection{Comparison with random matrix ensembles}

We compared the spectral properties of MEC dynamics against several random matrix ensembles (Table~\ref{tab:ensembles}). The MEC eigenvalue spectrum is best approximated by a shifted Wigner ensemble $A = W - cI$ with coupling strength $c \approx 0.35$ (distance $d = 0.51\sigma$ in $\alpha$), and by a generalized Ising model with zero-mean Gaussian couplings $J_{ij} \sim \mathcal{N}(0, \sigma^2/N)$ producing $\alpha \approx 0.076$ (distance $0.51\sigma$ from the MEC mean).

However, \textbf{eigenvector statistics} \cite{evers2008} reveal a more nuanced picture. Analysis of inverse participation ratio (IPR) \cite{wegner1980} and level spacing ratio distributions \cite{atas2013} across 15 MEC recordings shows that MEC eigenvectors differ from GOE at $3.47\sigma$ and from Ising criticality at $4.55\sigma$, but match a sparse random correlation ensemble (5\% density) within $0.86\sigma$ (Fig.~\ref{fig:eigenvectors}). This partial agreement is a finite-size coincidence: scaling analysis below shows the separation grows with $N$.

\begin{table}
\caption{Ensemble comparison. Distance $d$ (in standard deviations of the null) between MEC and each ensemble. Values $d < 1$ indicate indistinguishability.}
\label{tab:ensembles}
\begin{ruledtabular}
\begin{tabular}{lcc}
Ensemble & $d$ (eigenvalues) & $d$ (eigenvectors) \\
\hline
GOE & $>3$ & $3.47$ \\
Shifted Wigner ($g=0.35$) & $0.51$ & $3.45$ \\
Sparse ER ($p=0.05$) & $>3$ & $0.86$ \\
Ising critical & $0.51$ & $4.55$ \\
\end{tabular}
\end{ruledtabular}
\end{table}

\begin{figure}[t]
\centering
\includegraphics[width=\columnwidth]{figures/fig2_eigenvector_comparison.pdf}
\caption{\textbf{Eigenvector statistics.} (a) IPR distribution for MEC (blue) vs GOE (red) vs sparse ER $p=0.05$ (green). (b) Level spacing ratio distribution. (c) Ensemble distance summary.}
\label{fig:eigenvectors}
\end{figure}

\textbf{Scaling separation analysis:} To determine whether the MEC/sparse ensemble match persists at larger system sizes, we performed a subsampling scaling analysis \cite{fisher1971,binder1981}. Starting from 15 MEC recordings with $N \ge 100$, we subsampled neurons at sizes $N \in \{20, 30, 50, 80, 100\}$ and computed the distance $d(N)$ between MEC and sparse ensembles in eigenvector-statistic space. The distance grows as $d(N) \sim N^{2.2 \pm 0.1}$ (Fig.~\ref{fig:scaling}, inset), strongly rejecting the hypothesis that MEC converges to the sparse null at larger $N$. This growth is anomalous in the renormalization-group sense \cite{wilson1975,goldenfeld1992}: it implies a relevant operator whose effect amplifies with scale, rather than a finite-size transient. The result confirms that the observed correlation geometry reflects genuine structural organization rather than a finite-$N$ artifact.

\subsection{Finite-size scaling}

We analyzed the scaling of PR with system dimensionality $N$ (Fig.~\ref{fig:scaling}). MEC recordings show a saturating trend: PR approaches a plateau at large $N$ ($\text{PR} \approx 40$ for $N > 300$), consistent with exponential spectral truncation. Synthetic stabilized systems show power-law growth $\text{PR} \sim N^{0.125}$, consistent with heavy-tailed spectra.

\begin{figure}[t]
\centering
\includegraphics[width=\columnwidth]{figures/fig3_finite_size_scaling.pdf}
\caption{\textbf{Finite-size scaling of participation ratio.} PR vs dimensionality $N$ for MEC (blue circles) and synthetic systems (orange triangles). MEC shows saturating behavior; synthetic systems follow power-law growth. The black curve shows the prediction for exponential spectral truncation.}
\label{fig:scaling}
\end{figure}

This saturation is a key diagnostic: if PR($N$) asymptotes to a finite value, the spectrum is truncated at a characteristic scale rather than decaying as a true power law. This distinguishes the MEC regime from scale-invariant critical systems.

\subsection{Stability and generativity}

The product $\alpha \cdot \text{PR} \approx 2.3$ for MEC, satisfying the general inequality $\alpha \cdot \text{PR} \ge 2$ derived from the spectral bound. Systems near this bound trade spectral breadth (small $\alpha$) against effective dimensionality (large PR). The quantity $\text{PR} \cdot e^{\text{LC}/2} \approx 5.8 \pm 1.0$ is conserved across MEC sessions, indicating a coupling between dimensional richness and trajectory closure.

\subsection{Temporal interaction structure}

We recovered interaction operators from MEC activity using sparse inverse covariance estimation (graphical lasso) \cite{friedman2008,meinshausen2006}. The precision matrices show 60--89\% sparsity with spectral decay $\alpha \approx 0.33 \pm 0.11$ and PR $\approx 98 \pm 53$, confirming broad spectral support in the instantaneous correlation geometry \cite{dempster1972}. The precision matrix spectrum separates from sparse random ensembles at $3$--$4\sigma$, suggesting structure beyond simple sparsity \cite{ravikumar2011}.

Notably, temporal operators estimated via regularized vector autoregression show flat spectra ($\alpha \approx 0.01$, PR $\approx N$), indicating that lag-1 temporal dynamics are weak at the 20~ms timescale. The dominant structure is in instantaneous correlations rather than temporal propagation.

\section{Discussion}

\paragraph{Candidate stable broad-spectrum regime.}
The results are consistent with a stable broad-spectrum regime---provisionally referred to as Stably Sustained Critical Spectrum (SSCS)---characterized by $\alpha < 0.1$, PR $> 15$, and LC $< -5$. The scaling analysis shows that the MEC/sparse separation grows with system size ($d(N) \sim N^{2.2}$), strengthening the case for genuine structural organization beyond finite-size sparse correlations. However, this remains a candidate classification. The eigenvalue spectrum alone cannot distinguish MEC from the shifted Wigner ensemble or from Ising criticality \cite{hohenberg1977} (both at $d = 0.51\sigma$). The SSCS candidate differs from the self-organized criticality paradigm \cite{beggs2003,mora2011,munoz2018} in two respects: MEC dynamics are linearly stable (LC $< 0$, not marginally stable) and PR($N$) saturates at large $N$ (not power-law scaling).

\paragraph{Distinction from standard criticality.}
The SSCS candidate differs from conventional criticality in two respects. First, MEC dynamics are linearly stable (LC $< 0$), while critical systems are marginally stable. Second, the finite-size scaling of PR saturates, while critical systems typically exhibit power-law scaling. The spectral properties are better captured by a heterogeneous coupling model with $\alpha \approx 0.076$ than by an Ising model at $T_c$.

\paragraph{Distinction from collapse.}
Synthetic stabilized systems universally collapse to PR $\approx 1$, $\alpha > 0.5$, and LC $\approx 0$. The MEC regime is cleanly separated from collapse by all three metrics. The bound $\alpha \cdot \text{PR} \ge 2$ provides a rigorous separation criterion.

\paragraph{Scaling resolution.}
Scaling separation analysis resolves a previously open ambiguity. At fixed $N \approx 120$, MEC eigenvector statistics matched a sparse ensemble within $0.86\sigma$. This raised the possibility that the observed broad-spectrum organization was a finite-size artifact of sparse random correlations. However, subsampling analysis reveals that $d(N)$ grows as $N^{2.2}$, not toward zero. The MEC/sparse separation increases with system size, consistent with genuine structural organization.

\paragraph{Operator interpretation.}
The broad-spectrum organization resides in instantaneous correlation geometry, not in temporal propagation. Precision matrix spectra ($\alpha \approx 0.33$, PR $\approx 98$) show richer structure than lag-1 operators ($\alpha \approx 0.01$). The relevant object for understanding MEC organization is the instantaneous interaction network, not the temporal propagator. Higher-lag interactions and nonlinearities remain unexplored and may carry additional information.

\section{Methods}

\subsection{Data}

MEC recordings from 40 sessions of rodent spatial navigation (Calais, Goa, Kerala, Mumbai cohorts). Spike counts binned at 20~ms, Gaussian-smoothed ($\sigma = 40$~ms). Covariance matrices computed from $z$-scored activity.

\subsection{Spectral analysis}

Eigenvalues computed via SVD of the $z$-scored data matrix. Spectral decay $\alpha$ estimated by exponential fit $\log \lambda_k = -\alpha k + \text{const}$ to the leading half of the eigenvalue spectrum (see companion paper for bound derivation). IPR computed as $\text{IPR} = \sum_i |v_i|^4$ for normalized eigenvectors.

\subsection{Ensemble comparison}

Ensemble distances computed as $d = \sqrt{d_{\text{IPR}}^2 + d_r^2 + d_\alpha^2}$, where each component is the absolute difference normalized by the ensemble standard deviation.

\subsection{Operator recovery}

Graphical lasso with $\alpha = 0.1$ regularization, cross-validated when computationally feasible. VAR(1) estimates via ridge regression ($\ell_2$ penalty $\lambda = 1.0$). Dynamic mode decomposition via exact DMD.

\bibliography{../bibliography}

\end{document}
